\ssscp{Lengthened consonants in \tsc{Sumba}}{03-03-03-02}
\ili{As} discussed in \srf{03-03-01-01},
the languages of western Sumba have contrastive lengthened medial
consonants in certain words; e.g. \ili{Kodi} /potːo/ ‘k.o.
bamboo’ /poto/ ‘four’. One source of lengthening is historic schwa.
\ili{As} discussed in \srf{03-02-01} and \srf{03-03-01},
schwa usually triggers lengthening of a following consonant.
This is found in all the languages of Flores,
as well as \tsc{\ili{Havu}-Dhao},
and can be reconstructed to Proto-\ili{Bima}-Lembata.
The languages of western Sumba have lost schwa due to *ə > \it{a,
o, e} but usually retain the lengthening.
Examples are given in \trf{03-55}.
Phonetic transcriptions of the \ili{Havu} and \ili{Dhao} data
are given to highlight the (predictably) lengthened consonant

\begin{table}\tcp{\tsc{Sumba-Havu} C > Cː /*ə{\gap}}{03-55}
\begin{adjustbox}{width=\textwidth}
	\begin{threeparttable}\st{2pt}
	\begin{tabular}{llllllll}\lsptoprule
*gloss	&	CC		&	egg		&	turtle		&	full		&	day, sun		&	mouth		&	giant bamboo		\\
PMP	&	*əmbu\su{†}		&	*qatəluR		&	*pəñu		&	*ma-pənuq		&	*qaləjaw		&	*bəqbəq		&	*bətuŋ		\\	\midrule
\ili{Kodi}	&	{\hr}a\tbr{mbː}u\su{‡}	&	{\hqa}ta\tbr{lː}u	&	{\hr}ɣa\tbr{nː}u	&	{\hr}mba\tbr{nː}u	&	{\hqa}lo\tbr{ɗː}o\su{‡}	&	{\hr}ɣo\tbr{ɓː}a\su{‡}	&	{\hr}po\tbr{tː}o	\\
\ili{Weyewa}	&	\cg{(umbu)}		&	{\hqa}to\tbr{lː}u	&	{\hr}wo\tbr{nː}u	&	{\hr}mbo\tbr{nː}u\su{Δ}	&	{\hqa}lo\tbr{ɗː}o	&	{\hr}a\tbr{ɓː}a	&	{\hr}po\tbr{tː}o	\\
\ili{Loli}	&	\cg{(ubːu)}		&	{\hqa}to\tbr{lː}u	&			&			&	\hp{qa}‹lo\tbr{dd}o›	&	{\hr}aba		&	{\hr}po\tbr{tː}o	\\
\ili{Laboya}	&	\cg{(ubu)}		&	{\hqa}talu		&	{\hr}yanu		&	{\hr}bana		&	{\hqa}laɗo		&	{\hr}ɓaa		&			\\
\ili{Wanukaka}	&	\cg{(ubu)}		&	{\hqa}tilu		&	\cg{(tanoma)}		&	{\hr}bi\tbr{nː}u	&	{\hqa}ladu		&	{\hr}aa, aʔa		&	{\hr}pi\tbr{tː}iŋ	\\
\ili{Mamboru}	&	\cg{(umbuk-a)}	&	{\hqa}tilu		&	\cg{(tanoma)}		&	{\hr}mbinu		&	{\hqa}lɐɗu		&	{\hr}wɐɓa		&	{\hr}pɐtiŋ-u		\\
\ili{Anakalang}	&	\cg{(ubuk-u)}		&	{\hqa}tilu		&	\cg{(tanoma)}		&	{\hr}binu		&	\hp{qa}‹làdu›		&	‹yàba›		&	{\hr}pituŋ-u		\\
\ili{Lewa}	&	\cg{(umbuk-u)}	&	{\hqa}tilu		&	\cg{(tanoma)}		&	{\hr}mbinu		&			&	{\hr}yɐɓa		&			\\
\ili{Kambera}	&	\cg{(umbuk-u)}	&	{\hqa}tilu		&	\cg{(tanoma)}		&	{\hr}mbinu		&	{\hqa}loɗu		&	{\hr}yɐɓa		&	{\hr}pitiŋ-u		\\
\ili{Mangili}	&	\cg{(umbuk-u)}	&	{\hqa}tilu		&	\cg{(tanoma)}		&	{\hr}mbinu		&	{\hqa}lɐɗu		&	{\hr}yɐɓa		&	{\hr}pɐtiŋ-u		\\	\midrule
\ili{Havu}	&	\hp{[ˈʔ}əpu		&	{\hqa}dəlu		&	\hp{[ˈʔ}əɲu		&\cg{(tobo)}			&	{\hqa}loɗo		&	{\hr}uɓa\su{ǁ}	&			\\
	&	[ˈʔə\tbr{p}ːu]	&	\hp{qa}[ˈdə\tbr{lː}u]	&	[ˈʔə\tbr{ɲː}u]	&			&	\hp{qa}[ˈlo\tbr{ɗː}o]	&	[ˈuɓa]		&			\\
\ili{Dhao}	&	\hp{[ˈʔ}əpu		&	\cg{(kanaɖ͡ʐu)}	&	\hp{[ˈʔ}əɲu		&	{\hr}pənu		&	{\hqa}loɗo		&	{\hr}həɓa		&			\\
	&	[ˈʔə\tbr{pː}u]	&			&	[ˈʔə\tbr{ɲː}u]	&	[ˈpə\tbr{nː}u]		&	\hp{qa}[ˈlo\tbr{ɗː}o]	&	[ˈhə\tbr{ɓː}a]	&			\\
		\lspbottomrule
	\end{tabular}
		\begin{tablenotes}\tnsize
			\item[†]	{
								%Most words from Sumba listed under *əmpu here
								%are from the doublet *umpu without penultimate schwa.
								\ili{Kodi} \it{ambːu} is ‘grandparent’,
								\ili{Laboya} \it{ubu} ‘grandchild’, \ili{Wanukaka} \it{ubu} ‘grandparent,
								grandchild’. The other words here are from \citet[515]{Onvlee1984}
								in which \ili{Kambera} \it{umbuk-u} is ‘grandchild’.
								\ili{Havu} \it{əpu} is ‘grandparent,
								master’, \ili{Dhao} \it{əpu} is ‘grandparent,
								grandchild’.}
			\item[‡]	{Lengthening of prenasalised consonants in \ili{Kodi} is realised
								by lengthening the nasal component; thus /aᵐbu/ → [ˈʔɐmːbu] ‘grandparent’.
								Lengthening of an imploded stop is realised as a phonetic sequence
								of a glottal stop followed by an implosive: /loɗːo/ → [ˈlɔʔɗɔ]
								‘day, sun’, /ɣoɓːa/ → [ˈɣɔ̈ʔɓa] ‘mouth’.
								Phonetic data on the realisation of lengthened consonants in
								other languages of Sumba is not available.}
			\item[Δ]	{\citet[74]{NgongoBili2020} has \ili{Weyewa} \it{bonːu} ‘full’.}
			\item[{ǁ}]	{The \ili{Dimu} dialect of \ili{Havu} has \it{vuɓa} ‘mouth’.}
		\end{tablenotes}
	\end{threeparttable}
\end{adjustbox}
\end{table}

While most cases of historical schwa trigger lengthening in the
languages of western Sumba,
there are also occasional exceptions.
Three exceptions in \ili{Kodi} are *təlu > \it{talu} ‘three’ (contrasting
with *qatəluR > \it{talːu} ‘egg’ in \trf{03-55}),
*təbəl > \it{tembe} ‘thick’,
and *dəŋəR > \it{roŋo} ‘hear’ (but \ili{Weyewa} \it{reŋːe} ‘hear’).

%Another source of lengthening in western Sumba may be shared with \tsc{\ili{Havu}-Dhao}.
%This concerns words with a penultimate high vowel,
%*i or *u, and final low vowel, *a or *ə.
%\tsc{\ili{Havu}-Dhao} have metathesis of the vowels in such words along
%with *a > \it{ə}.
%(Schwa in \ili{Havu} and \ili{Dhao} bears stress
%and triggers lengthening of a following consonant.) In the same
%environment, the languages of western Sumba,
%in particular \ili{Kodi}, \ili{Weyewa},
%and \ili{Loli}, usually have lengthening of the medial consonant.
Lengthened consonants also develop in western Sumba
for words with a penultimate high vowel
(*i or *u) and final non-high vowel (*a or *ə).
That is, in the environment *V\tsc{[+high]}CV\tsc{[-low]}{\#} > VCːV{\#}.
Note that this is the same environment in which vowel metathesis
occurs in \tsc{\ili{Havu}-Dhao} (see \trf{02-10b} \prf{tab:02-10b}).
%\ili{Wanukaka} has sporadic lengthening in the same environment,
%while other languages of Sumba for which data is available have no lengthening.
Examples are given in \trf{03-56}.
%Phonetic transcriptions of the \ili{Havu}
%and \ili{Dhao} reflexes are given to highlight the stress bearing schwa
%and lengthening of the following consonant.

\begin{table}\tcp{\tsc{Sumba-Havu} *V\tsc{[+high]}CV\tsc{[+low]}{\#}}{03-56}
\begin{adjustbox}{width=\textwidth}
	\begin{threeparttable}\st{2pt}
	\begin{tabular}{llllllll@{}l}\lsptoprule
*gloss	&	house	&	moon	&	clay pot		&	vomit	&	oil	&	see	&	five	&	woman\\
PMP	&	*Ru\tbr{m}aq	&	*bu\tbr{l}an	&	*ku\tbr{d}ən		&	*um-u\tbr{t}aq	&	*mi\tbr{ñ}ak	&	*ki\tbr{t}a	&	*li\tbr{m}a	&	**bi\tbr{n}ai\su{‡}\\	\midrule
\ili{Kodi}	&	\hp{*R}u\tbr{mː}a	&	{\hr}wu\tbr{lː}a	&	{\hr}ɣu\tbr{rː}a\su{†}	&	{\hr}mu\tbr{tː}a	&	{\hr}mi\tbr{ɲː}a	&	{\hr}yi\tbr{cː}e	&	{\hr}li\tbr{mː}a	&	{\hr}kawi\tbr{ɲː}a\\
Wey.	&	\hp{*R}u\tbr{mː}a	&	{\hr}wu\tbr{lː}a	&	{\hr}wi\tbr{rː}o\su{†}	&	‹mu\tbr{t}a›	&	\cg{(leŋi)}	&	{\hr}eta	&	{\hr}li\tbr{mː}a	&	{\hr}mawi\tbr{nː}e\\
\ili{Loli}	&	\hp{*R}u\tbr{mː}a	&	{\hr}wu\tbr{lː}a	&	{\hr}wu\tbr{rː}o		&	{\hr}mu\tbr{tː}a	&		&	{\hr}ieta	&	{\hr}li\tbr{mː}a	&	\\
\ili{Laboya}	&	\hp{*R}uma	&	{\hr}wula	&			&	{\hr}muta	&		&	{\hr}eːta	&	{\hr}lima	&	{\hr}aŋu wine\\
Wanuk.	&	\hp{*R}uma	&	{\hr}wula\cg{(ŋ-u)}	&	{\hr}a\tbr{rː}u\cg{(ŋ)}	&	{\hr}muta	&	\cg{(leŋih)}	&	{\hr}i\tbr{tː}a	&	{\hr}lima, li\tbr{mː}a&	{\hr}mawini\\
Mamb.	&	\hp{*R}uma	&	{\hr}wula	&	\cg{(kaseri)}		&	{\hr}muta	&	{\hr}miɲa	&	{\hr}ita	&	{\hr}lima	&	{\hr}kawini\\
Anakal.	&	\hp{*R}uma	&	{\hr}wula	&	{\hr}ɐruŋ-u		&	{\hr}muta	&	{\hr}miɲa	&	{\hr}ita	&	{\hr}lima	&	{\hr}mawini\\
\ili{Lewa}	&	\hp{*R}uma	&	{\hr}wulaŋ-u	&	{\hr}ɐruŋ-u		&	{\hr}muta	&	{\hr}mina	&	{\hr}ita	&	{\hr}lima	&	{\hr}wowinu\\
Kamb.	&	\hp{*R}uma	&	{\hr}wulaŋ-u	&	{\hr}wuruŋ-u		&	{\hr}muta	&	{\hr}mina	&	{\hr}ita	&	{\hr}lima	&	{\hr}kawini\\
\ili{Mangili}	&	\hp{*R}uma	&	{\hr}wulaŋ-u	&	{\hr}wuraŋ-u		&	{\hr}muta	&	{\hr}mina	&	{\hr}ita	&	{\hr}lima	&	{\hr}kawini\\\midrule
\ili{Havu}	&	\hp{*R}əmu	&	{\hr}vəru	&	{\hr}əru		&	{\hr}mədu	&	{\hr}məɲi	&	{\hr}ŋədi	&	{\hr}ləmi	&	{\hr}bəni\\
	&	[ˈʔə\tbr{mː}u]	&	[ˈvə\tbr{rː}u]	&	[ˈə\tbr{rː}u]		&	[ˈmə\tbr{dː}u]	&	[ˈmə\tbr{ɲː}i]	&	[ˈŋə\tbr{dː}i]	&	[ˈlə\tbr{mː}i]	&	[ˈbə\tbr{nː}i]\\
\ili{Dhao}	&	\hp{*R}əmu	&	{\hr}həru	&	{\hr}əru		&	{\hr}məɖ͡ʐu	&	{\hr}məɲi	&	-əɖ͡ʐi	&	{\hr}ləmi	&	{\hr}b͡βəni\\
	&	[ˈʔə\tbr{mː}u]	&	[ˈhə\tbr{rː}u]	&	[ˈə\tbr{rː}u]		&	[ˈmə\tbr{dːʐ}u]&	[ˈmə\tbr{ɲː}i]	&	[ˈə\tbr{dːʐ}i]	&	[ˈlə\tbr{mː}i]	&	[ˈbβə\tbr{nː}i]\\
		\lspbottomrule
	\end{tabular}
		\begin{tablenotes}\tnsize
			\item[†]	{\citet[548]{Onvlee1984} gives \ili{Kodi} ‹ghurongo› and \ili{Weyewa} ‹wuro›.
								\ili{Weyewa} \it{wirːo} is from \citep[69]{Bulu2018}.}
			\item[‡]	{\ili{Kodi} \it{kawiɲːa} is ‘widow’.
								\citet{Mitchell2008b} gives \ili{Loli} ‹mawine ?› for ‘woman’.
								\ili{Laboya} \it{aŋu wine} is ‘woman’s younger sister’.}
		\end{tablenotes}
	\end{threeparttable}
\end{adjustbox}
\end{table}

While it is usual for the medial consonant in *V\tsc{[+high]}CV\tsc{[-high]}{\#}
sequences to undergo lengthening in western Sumba,
there are also a small number of examples where this does not occur.
Two such examples are: *ina > \ili{Kodi} \it{iɲa} ‘mother’,
\ili{Havu} \it{ina} ‘mother, female’, \ili{Dhao} \it{ina} ‘maternal aunt’; as well
as *pija > \ili{Kodi} \it{pira} ‘when’ (but \ili{Weyewa} \it{pirːa},
\ili{Havu}, \ili{Dhao} \it{pəri} ‘how many’).

\ili{As} already discussed in \srf{03-03-02-HD}, and highlighted
by the phonetic transcriptions above, penultimate schwa
in \tsc{\ili{Havu}-Dhao} triggers automatic lengthening of the following consonant.
This applies to both retentions of schwa from PMP, such as
qatəluR > \ili{Havu} \it{dəlu} [ˈdə\tbr{lː}u] ‘egg’
(see \trf{02-09a} \prf{tab:02-09a} for more examples),
as well as instances of schwa that became penultimate due to metathesis,
such as *Rumaq > \ili{Havu} \it{əmu} [ˈʔə\tbr{mː}u].
In the latter case, the lengthening is an innovation,
which may be due to analogy with
other instances of penultimate schwa.

Although the lengthening of medial consonants in \ili{Havu}
and \ili{Dhao} in examples such as those in \trf{03-56}
can be analysed synchronically as conditioned by the
previous schwa it seems unlikely to be a
coincidence that the languages of western Sumba have lengthening
in the same (historical) environment,
that is words with the (historic) shape *V\tsc{[+high]}CV\tsc{[+low]}.
Given that the languages of Sumba do not have
metathesis or the development of schwa in these
words, the lengthening here cannot be analysed as
an innovation dependent on another innovation,
as it can be in \tsc{\ili{Havu}-Dhao}.

This raises the possibility that the lengthening
in words with this shape originally arose in \tsc{\ili{Havu}-Dhao}
before the vowel metathesis.
If so, lengthening of medial consonants between a high vowel
and low vowel could be a development shared between \tsc{Sumba}
and \tsc{\ili{Havu}-Dhao}, with vowel metathesis a
subsequent development in \tsc{\ili{Havu}-Dhao}.

\begin{table}\tcp{\tsc{Sumba-Havu} \tsc{[+high]}CV\tsc{[+high]} and *aCV}{03-57}
\begin{adjustbox}{width=\textwidth}
	\st{2pt}\begin{tabular}{lllllllll}\lsptoprule
*gloss	&	seed	&	hair	&	black	&	child	&	chicken	&	rice plant	&	drink	&	fire\\
PMP	&	*binəhiq	&	*bulu	&	*ma-qitəm	&	*anak	&	*manuk	&	*pajay	&	*inum	&	*hapuy\\	\midrule
\ili{Kodi}	&	{\hr}wini	&	{\hr}wulu	&	{\hr}mete	&	{\hr}ana	&	{\hr}manu	&	{\hr}pari	&	{\hr}inu	&	{\hr}api\\
\ili{Weyewa}	&	{\hr}wi\tbr{nː}i	&	{\hr}wu\tbr{lː}u	&	{\hr}me\tbr{tː}e	&	{\hr}ana	&	{\hr}manu	&	{\hr}pare	&	{\hr}enu	&	{\hr}api\\
\ili{Loli}	&	{\hr}wi\tbr{nː}i	&	{\hr}wu\tbr{lː}u	&	{\hr}me\tbr{tː}e	&	{\hr}ana	&	{\hr}manu	&	{\hr}pare	&	{\hr}ienu	&	{\hr}api\\
\ili{Laboya}	&		&	{\hr}wulu	&	{\hr}mati	&	{\hr}a:na	&	{\hr}maːnu	&	{\hr}paːre	&	{\hr}eːnu	&	\cg{(lati)}\\
\ili{Wanukaka}	&	{\hr}wini	&	{\hr}wulu	&	{\hr}metiŋ-u	&	{\hr}ana	&	{\hr}manu	&	{\hr}pari	&	{\hr}inu	&	{\hr}api\\
\ili{Mamboru}	&	{\hr}wini	&	{\hr}wulu	&	\cg{(mɐli)}	&	{\hr}ana	&	{\hr}manu	&	{\hr}pari	&	{\hr}unu	&	{\hr}api\\
\ili{Anakalang}	&	{\hr}wini	&	{\hr}wulu	&	{\hr}metiŋ-u	&	{\hr}ana	&	{\hr}manu	&	{\hr}pari	&	{\hr}inu	&	{\hr}api\\
\ili{Lewa}	&	{\hr}winu	&	{\hr}wulu	&	{\hr}metuŋ-u	&	{\hr}ana	&	{\hr}manu	&	{\hr}pari	&	{\hr}uniŋ-u	&	{\hr}apu\\
\ili{Kambera}	&	{\hr}wini	&	{\hr}wulu	&	{\hr}mitiŋ-u	&	{\hr}ana	&	{\hr}manu	&	{\hr}pari	&	{\hr}unuŋ-u	&	{\hr}epi\\
\ili{Mangili}	&	{\hr}wini	&	{\hr}wulu	&	{\hr}mitiŋ-u	&	{\hr}ana	&	{\hr}manu	&	{\hr}pari	&	{\hr}unuŋ-u	&	{\hr}api\\\midrule
\ili{Havu}	&	{\hr}vini	&	\cg{(rəu)}	&	{\hr}mədi	&	{\hr}ana	&	{\hr}manu	&	\hp{*p}are	&	{\hr}ŋinu	&	{\hr}ai\\
\ili{Dhao}	&	{\hr}hini	&	\cg{(rəu)}	&	{\hr}məɖ͡ʐi	&	{\hr}ana	&	{\hr}manu	&	\hp{*p}are	&	{\hr}-inu	&	{\hr}ai\\
		\lspbottomrule
	\end{tabular}
\end{adjustbox}
\end{table}

If this is the case, then it may help shed some light
on how exactly the \tsc{\ili{Havu}-Dhao} vowel metathesis
came about.\footnote{
		\citet[220--229]{Blust2012a} considers possible
		origins of the vowel metathesis in \tsc{\ili{Havu}-Dhao},
		but finds none of them entirely satisfactory.
		Blust does not consider whether consonant lengthening played a
		role in the process.
		Although the data from Sumba indicates this may be relevant factor,
		the way in which it may be relevant remains to be worked out.}
Thus, for instance, \citet[48]{Lovestrand2021} observes that
vowels before lengthened consonants in \ili{Kodi} are shorter than those before single consonants.
This may indicate that the schwa in \tsc{\ili{Havu}-Dhao} is a result
of vowel reduction *V\tsc{[+high]} > \it{ə} /{\gap}CːV{\#}.
Nonetheless, the source of the final high vowel would still require explanation.
More detailed data from more of the languages of western Sumba -- in
particular acoustic phonetic studies of the vowels -- may shed
some light on the matter.

For the time being, we do not treat the lengthening
of medial consonants in western Sumba and \tsc{\ili{Havu}-Dhao}
in words of the shape *V\tsc{[+high]}CV\tsc{[+low]} as
a shared innovation, mainly because the \tsc{\ili{Havu}-Dhao}
lengthening is analysable as being triggered by the vowel metathesis.
However, the possibility that the two cases of lengthening are
somehow linked deserves further investigation when more
date from languages of western Sumba comes to light.

%The shared patterns regarding *V\tsc{[+high]}CV\tsc{[-high]}{\#} words in
%western Sumba and \tsc{\ili{Havu}-Dhao} may also point to \tsc{\ili{Havu}-Dhao} being more
%closely related to the languages of western Sumba than they are
%to the those of eastern Sumba.
%However, it is also possible that consonant lengthening developed
%in Proto-Sumba-\ili{Havu}, but was lost in the languages of eastern Sumba.
%Indeed, this appears to be what has happened in \ili{Laboya} which
%has lost distinctive consonant lengthening,
%but developed distinctive vowel length (see \srf{03-03-01},
%\citealp[48]{Lovestrand2021} and \citealp[18]{Verdizade2019a}).

\ili{Weyewa} and \ili{Loli} also have lengthening of consonants between two
high vowels, as well as in some other cases involving other vowels.
Lengthening does not usually occur in other languages of Sumba
apart from the cases discussed above, but exceptions occur.
Four examples from \ili{Kodi} which have lengthening are: *i-kami >
\it{ya\tbr{mː}a} ‘\tsc{1pl.excl}’, *i-kamuyu > \it{ye\tbr{mː}i} ‘\tsc{2pl}’,
*masik > \it{ma\tbr{hː}i} ‘salt’,
and *wahiyəR > \it{we\tbr{yː}o} ‘water’.
Examples of the reflexes of words which are not *V\tsc{[+high]}CV\tsc{[+low]}{\#}
or *əCV{\#} are given in \trf{03-57}.
